\section{Introduction}

Modern distributed systems face growing data volumes and tighter response-time requirements as the number of users and services scales \cite{kleppmann2017}. To support concurrent read and write workloads, the Command Query Responsibility Segregation (CQRS) pattern has become a widely adopted approach that separates write (\textit{command}) and read (\textit{query}) paths, allowing each side to be optimized independently \cite{fowler2011,munonye2023}.

The separation, however, introduces a new burden: the read model must be kept synchronized with the write model. In typical implementations, write-side changes are emitted as events and propagated to the read side through a message broker, often using event sourcing \cite{morris2023}. Synchronization becomes asynchronous and exposes the system to \textit{eventual consistency}, where the read model temporarily lags behind the write model \cite{vogels2009,bailis2012}.

Two parameters dominate this asynchronous path: \textit{batch size} and \textit{poll interval}. Under high load, larger batches amortize per-message overhead such as network transfers, compression, and log writes, thus improving throughput. Larger batches, however, raise end-to-end latency because consumers wait longer to accumulate events, and large batches concentrate load on the read database, causing CPU, memory, and I/O spikes. Small batches reduce latency but inflate per-message overhead and increase commit frequency \cite{wu2019,rahman2010}. Setting these parameters statically is therefore unable to track the workload over time.

Prior work has explored dynamic batching. Wu et al. \cite{wu2020} proposed a reactive batching strategy in Apache Kafka that adapts batch size based on latency and throughput metrics. Das et al. \cite{das2018} introduced adaptive micro-batching in Spark Structured Streaming. These approaches improve performance, but most rely on heuristics or threshold rules tied to a single observed metric. They do not provide a formal framework for jointly balancing multiple competing variables such as backlog, latency, CPU, and memory utilization, and they offer no explicit notion of optimality.

This paper formulates CQRS synchronization as a discrete-time linear dynamic system and applies a Linear Quadratic Regulator (LQR) to compute control actions for the batch size and poll interval. The state vector aggregates backlog, CPU, memory, and I/O utilization, while the cost function captures the trade-off between synchronization quality and resource usage through the weighting matrices $Q$ and $R$. The contributions of this work are:

\begin{enumerate}
    \item A state-space formulation of CQRS event-sink synchronization as a MIMO control problem, including a measurement procedure for end-to-end latency via manual offset commit.
    \item An LQR controller whose gain matrix is derived from an identified linear model of the sink dynamics and whose cost function explicitly encodes the trade-off between backlog reduction and resource pressure.
    \item An empirical comparison against static, rule-based, and PID baselines on five workload patterns (step, ramp, impulse, periodic step, low step) and a sensitivity analysis over three weighting configurations using \textit{normalized regret}.
\end{enumerate}

The remainder of this paper is organized as follows. Section~2 reviews CQRS, stream processing, and control approaches. Section~3 presents the system model and the LQR-based controller. Section~4 describes the experimental setup, load scenarios, and evaluation metrics. Section~5 reports and discusses the results. Section~6 concludes.
